\chapter*{Conclusion Générale et Perspectives}
\addcontentsline{toc}{chapter}{Conclusion Générale et Perspectives}

Le projet Fraudly représente une exploration ambitieuse et réussie des formidables opportunités offertes par la convergence de l'ingénierie logicielle distribuée et de l'intelligence artificielle agentique dans le secteur de l'éducation. En remplaçant les systèmes d'apprentissage classiques par un réseau de microservices robustes, nous avons posé les fondations technologiques nécessaires à la scalabilité d'une institution éducative moderne.

L'intégration minutieuse des grands modèles de langage a démontré qu'il est désormais possible d'offrir à chaque étudiant une expérience de tutorat sur mesure, capable de s'adapter à ses lacunes spécifiques. Parallèlement, le développement d'un module de proctoring multimodal complexe, alliant analyse comportementale en temps réel et vision par ordinateur avancée, apporte une réponse crédible aux menaces pesant sur l'intégrité de l'évaluation à distance.

L'aboutissement de cette phase du projet ouvre des perspectives de développement fascinantes. Sur le court terme, l'optimisation des modèles d'intelligence artificielle pour réduire davantage l'empreinte computationnelle permettra un traitement plus économique des flux de données. À plus long terme, l'intégration de capacités multimodales au sein des modèles d'évaluation permettrait la correction automatisée de schémas d'ingénierie ou de démonstrations mathématiques complexes.
